% ============================================================
% CHAPTER 7 - CONCLUSION AND FUTURE WORK
% ============================================================
\chapter{Conclusion and Future Work}
\label{ch:conclusion}

\section{Introduction}

This chapter concludes the dissertation. Section~\ref{sec:conc-summary}
summarises the work; Section~\ref{sec:conc-rq} answers each research question
against the evidence presented; Section~\ref{sec:conc-contributions} restates
the contributions; Section~\ref{sec:conc-implications} discusses the
implications; Section~\ref{sec:conc-reflection} reflects critically on the
process; Section~\ref{sec:future} sets out future work; and
Section~\ref{sec:conc-closing} closes.

\section{Summary of the Work}
\label{sec:conc-summary}

This dissertation set out to give a person with dysarthric or otherwise unclear
speech the kind of feedback a therapist gives: not ``that word was unclear'' but
``that \emph{sound} was unclear''. The review in Chapter~\ref{ch:literature}
established that the two mature technical traditions adjacent to this goal ---
mispronunciation detection for second-language learners, and recognition and
severity estimation for disordered speech --- each supply one half of what is
required and neither supplies both, and that the binding constraint is the
absence of clinician-annotated per-phoneme ground truth.

Chapter~\ref{ch:methodology} addressed that constraint by deriving labels from
the paired structure of the UASpeech corpus. Both recordings of a pair are
force-aligned to the same canonical phoneme sequence obtained from the control
transcript, so that phoneme $i$ of the impaired production and phoneme $i$ of
the healthy production are two attempts at the same intended sound and their
log-mel distance is meaningful. A first labelling strategy that thresholded the
pooled deviation distribution at a fixed percentile was implemented, found to
mislabel healthy speech by construction, and replaced by the control-baseline
strategy in which control instances are correct by definition and the error
cutoff is estimated over the dysarthric distribution alone.

Chapters~\ref{ch:analysis-design} and \ref{ch:implementation} designed and built
the system around the model: a CNN--BiLSTM classifier of 415{,}810 parameters, a
FastAPI inference service that reproduces the training feature pipeline exactly
and adds a quality gate and a target-word verifier, and a React progressive web
application that records, submits and visualises.

Chapter~\ref{ch:results} evaluated the result. Five architectures were compared;
CNN--BiLSTM was selected on error-class recall and $F_1$; the selected model's
behaviour was analysed in detail, including a frank account of its low precision
under natural class imbalance; the per-phoneme error profile was shown to
reproduce four independently documented clinical patterns; and the
control-versus-dysarthric separation --- the property the labelling redesign was
built to deliver --- was confirmed.

\section{Answers to the Research Questions}
\label{sec:conc-rq}

\subsection*{RQ1 --- How can phoneme-level labels be derived without clinician
annotation, and why does global percentile labelling fail?}

Labels can be derived by \emph{reference comparison}: aligning a healthy and an
impaired production of the same word to the same canonical phoneme sequence and
measuring the per-phoneme log-mel deviation between corresponding segments
(Equation~\eqref{eq:distance}). Global percentile labelling then fails for a
structural reason. Applying a percentile cutoff to the pooled distribution makes
the criterion ``is among the most deviant $\rho\%$ of everything'' rather than
``deviates from healthy speech''. A fixed fraction of healthy control phonemes
necessarily exceeds that cutoff and is labelled erroneous. The resulting label
noise is not random but systematic, and it is concentrated precisely on the
examples that ought to define the correct class, so the model learns the base
rate rather than the phenomenon
(Sections~\ref{sec:method1} and~\ref{sec:res-separation}).

\subsection*{RQ2 --- Does control-baseline labelling improve separation between
clear and impaired speech?}

Yes, and the improvement is qualitative rather than incremental. Under Method~1
healthy control recordings were returned with a large fraction of their phonemes
flagged; under Method~2 they return few or none, while dysarthric recordings of
the same words return substantially more
(Sections~\ref{sec:res-separation} and~\ref{sec:qualitative}). The mechanism is
explicit in Equation~\eqref{eq:method2}: all 49{,}153 control instances are
labelled correct unconditionally and retained in training as anchors, and the
cutoff is estimated over the dysarthric distribution alone, so no control
instance can enter the error class by any route. Since the architecture,
features and training protocol were essentially unchanged between the two
methods, the change in behaviour is attributable to the labelling strategy
alone.

\subsection*{RQ3 --- Which architecture performs best, and by what criterion?}

\textbf{CNN--BiLSTM}, selected on error-class recall (76.64\%) with error-class
$F_1$ (74.04\%) as tie-breaker; both are the best values in the study
(Table~\ref{tab:arch-comparison}). The criterion matters as much as the winner.
Accuracy does not discriminate --- the top three models span 0.27 percentage
points, which is inside run-to-run variation --- and accuracy is in any case
misleading under 85\% class imbalance, where a constant predictor scores 85\%.
Recall on the error class is the operationally correct criterion because the
costs are asymmetric: a false negative silently withholds the feedback the
system exists to provide, whereas a false positive costs one unnecessary
repetition. The ablation is informative in its own right: removing the recurrent
stage entirely costs 8.2 points of recall and 4.2 points of $F_1$, confirming
that temporal modelling of the phoneme's trajectory contributes measurably.

\subsection*{RQ4 --- Which phonemes are most frequently flagged, and is the
distribution clinically plausible?}

The highest predicted error rates fall on the alveolar flap \ph{ɾ} (50.84\%),
the rhotic schwa \ph{ɚ} (49.71\%), the velar nasal \ph{ŋ} (46.44\%) and the
reduced vowel \ph{ə} (45.58\%) (Table~\ref{tab:per-phoneme}). The distribution
is clinically plausible on four independent counts: rapid ballistic gestures
rank highest, rhotics and centralised vowels rank high (consistent with vowel
space compression), fricatives are over-represented relative to their frequency
(consistent with impaired airflow control), and nasals appear where
velopharyngeal control is implicated. That an ordering derived from a purely
data-driven acoustic distance measure reproduces these documented patterns is
evidence that the model has learned articulatory difficulty rather than a
dataset artefact (Section~\ref{sec:per-phoneme}).

\subsection*{RQ5 --- Can it be deployed accessibly on commodity hardware?}

Yes. The system runs as a single Docker image on a 2-vCPU CPU-only host with no
GPU, serving both the API and the compiled web client from one origin. Cold
start is 10--30\,s, paid once at process start and absorbed by the warm-up hook;
warm single-word analysis completes in a few seconds, dominated by the wav2vec2
alignment pass rather than by the classifier. The client is a progressive web
application installable on a phone home screen
(Sections~\ref{sec:impl-deploy} and~\ref{sec:res-system}).

\section{Contributions Restated}
\label{sec:conc-contributions}

\begin{enumerate}[label=\textbf{C\arabic*.},leftmargin=1.4cm]
  \item A documented diagnosis of the global-percentile labelling failure for
        paired impaired/healthy speech corpora, with the mechanism made
        explicit.
  \item The control-baseline labelling strategy of
        Equation~\eqref{eq:method2}, in which healthy productions are correct by
        definition and retained as anchors, and the error cutoff is estimated
        over the impaired sub-population only.
  \item The matched-pair alignment procedure that aligns both recordings to the
        same control-derived canonical sequence, which is what makes a
        per-phoneme distance well defined.
  \item A controlled five-architecture comparison with the selection criterion
        stated in advance and justified by the asymmetric costs of a screening
        application.
  \item A per-phoneme error profile validated against four independently
        documented clinical patterns.
  \item \systemName: a deployed reference implementation including a
        train--serve parity test, an audio quality gate and a target-word
        verifier --- production concerns absent from the research pipeline.
\end{enumerate}

\section{Implications}
\label{sec:conc-implications}

\textbf{For research on weakly supervised speech assessment.} Where ground truth
must be derived rather than annotated, the derivation deserves at least the
scrutiny normally given to the model. This work found a 3.75-point accuracy
spread across five architectures and a qualitative behavioural change from one
labelling decision. Papers reporting derived labels should state the derivation
rule explicitly and should test the resulting model on data that the rule
implicitly defines as negative --- here, healthy speech --- because that is where
a labelling pathology becomes visible.

\textbf{For clinical practice.} The result is not a diagnostic instrument and
must not be presented as one. What it plausibly supports is between-session
practice: a client can obtain per-sound feedback during unsupervised repetition,
and a clinician can review which sounds are consistently flagged. The interface
decisions of Section~\ref{sec:ui-design} --- exposing the probability, exposing
the threshold, never presenting a diagnosis --- are how that framing is enforced
in the artefact rather than only in the documentation.

\textbf{For accessible speech technology.} The system runs on free-tier CPU
hardware in a browser. The engineering choices that made this possible --- a
small classifier, client-side format conversion, CPU-only dependency wheels,
baked-in normalisation constants --- are transferable to other assistive speech
applications, where deployment cost is frequently the barrier between a
research result and a usable tool.

\section{Critical Reflection}
\label{sec:conc-reflection}

Three judgements would be made differently with hindsight.

\textbf{The validation and test partitions should never have coincided}
(limitation L2). Supplying the held-out partition to Keras as the validation set
was expedient and is a genuine methodological defect; the reported 71.99\% is
mildly optimistic as a result. A three-way split costs almost nothing and would
have removed an avoidable caveat from every number in Chapter~\ref{ch:results}.

\textbf{Each architecture should have been trained more than once.} With a
single run per candidate and a 0.27-point accuracy spread across the top three,
the comparison cannot support fine distinctions. The recall differences are
large enough to be believable, but the study would be stronger with a mean and a
standard deviation over several seeds.

\textbf{A small clinician-annotated validation set would have been worth its
cost.} Even one or two hundred phoneme instances judged by a speech-language
pathologist would convert the central caveat of this work --- that all metrics
measure agreement with a derived proxy --- from an unbounded limitation into a
measured one. This is the highest-value item of future work per unit of effort.

Against these, the decision to treat the labelling failure as the central
research problem rather than as a bug to be patched was correct, and it is what
gives the dissertation a result worth reporting.

\section{Future Work}
\label{sec:future}

\subsection{Immediate}

\begin{enumerate}[leftmargin=*]
  \item \textbf{Three-way split.} Separate validation and test partitions, and
        re-report all metrics on genuinely held-out data (addresses L2).
  \item \textbf{Repeated runs.} Train each architecture over several seeds and
        report means and standard deviations (addresses L6).
  \item \textbf{Threshold characterisation.} Compute the full precision--recall
        and ROC curves over the threshold range and publish the operating
        characteristic, so that the sensitivity slider is annotated with its
        actual consequences rather than a bare number.
  \item \textbf{Act on alignment confidence.} Treat a very low alignment score
        as a candidate omission and report it as a distinct verdict from a
        distortion (addresses L5).
\end{enumerate}

\subsection{Medium term}

\begin{enumerate}[leftmargin=*, resume]
  \item \textbf{Clinician-annotated validation.} Have speech-language
        pathologists annotate a sample of phoneme instances and measure the
        model's agreement with expert judgement, along with inter-annotator
        agreement as an upper bound (addresses L1).
  \item \textbf{Error typology.} Extend the binary verdict to distinguish
        substitution, distortion and omission, which is what makes feedback
        prescriptive rather than merely diagnostic.
  \item \textbf{Speaker adaptation.} Calibrate the decision threshold, or
        fine-tune the final layers, on a small enrolment sample from an
        individual user, following the personalisation results reported for
        disordered-speech recognition \cite{macdonald2021disordered}.
  \item \textbf{Connected speech.} Extend from isolated words to phrases and
        sentences, which introduces coarticulation, and to prosodic assessment.
  \item \textbf{Domain robustness.} Quantify and reduce the gap between
        studio-quality training recordings and consumer microphones through
        augmentation with noise, reverberation and codec artefacts
        (addresses L7).
\end{enumerate}

\subsection{Longer term}

\begin{enumerate}[leftmargin=*, resume]
  \item \textbf{Longitudinal tracking.} Persist analyses per user, with consent,
        and surface progress over time on specific phoneme targets --- the single
        most requested capability in therapy contexts and the one that would
        most change the tool's role from screening to intervention support.
  \item \textbf{Prescriptive feedback.} Map a flagged phoneme onto a targeted
        exercise and a plain-language articulatory instruction, closing the loop
        from detection to practice.
  \item \textbf{Multilingual extension.} \texttt{espeak-ng} and the wav2vec2
        aligner are both multilingual; the classifier is the language-specific
        component, and cross-lingual transfer of the control-baseline labelling
        strategy to a paired corpus in another language is a direct test of the
        method's generality.
  \item \textbf{Clinical study.} A trial with speakers with dysarthria measuring
        usability, comprehensibility and any effect on practice volume or
        outcome, which is the only way to establish whether the tool delivers
        the benefit it is designed for (addresses L8).
\end{enumerate}

\section{Closing Remarks}
\label{sec:conc-closing}

The most instructive outcome of this project was not the model. It was the
discovery that a labelling rule which is statistically unobjectionable ---
``call the most deviant thirty per cent errors'' --- can be quietly and
completely wrong, and that the wrongness is invisible in the training metrics
and obvious the moment the system is shown a healthy speaker. The fix required
no new architecture and no additional data: it required deciding that
\emph{correct} means \emph{like a healthy speaker producing this sound}, and
then enforcing that definition in the code.

\systemName\ is a research prototype and its limitations are stated plainly in
Section~\ref{sec:limitations}. It does not diagnose, it has not been validated
against clinicians, and roughly two of every three sounds it flags may not be
faults. What it does do is take a recording of a single word and say which
sound, at which moment, deviated from clear speech --- which is a thing a person
practising alone at home currently has no way to find out. Making that
information available, on a phone, at no cost, seems a worthwhile place for the
work to have arrived.
